\subsubsection{Effectiveness on Logic Queries}

As shown in \tref{tab:comparison_lq-dataloc}, \tooldataloc consistently outperforms all baselines across metrics and localization granularities on \dataset. At the file level, \tooldataloc with Qwen3-Max achieves the best Precision (PRE) and Success Rate (SR), reaching 73.35\% and 68.44\%, respectively. It also attains a Perfect Localization Rate (PLR) of 48.44\%, whereas all baselines remain near zero. Even at the finest function-level granularity, \tooldataloc still maintains a PLR of 38.27\%, while all baselines drop to 0\%. These results highlight the advantage of \tooldataloc in capturing code structure and reasoning capabilities in assisting precise localization. \tooldataloc maintains strong performance across both evaluated foundation models. Switching the underlying model from Claude-3.5 to Qwen3-Max yields gains across most localization metrics at all three granularities, although hit rate decreases slightly.

\begin{landscape}
\begin{table}[p]
    \centering
    \footnotesize
    \setlength{\tabcolsep}{2pt}
    \caption{Evaluation results on \dataset.}
    \label{tab:comparison_lq-dataloc}

    \begin{tabularx}{\linewidth}{@{} l l *{18}{C} @{}}
        \toprule
        \multirow{2}{*}{\textbf{Method}} & \multirow{2}{*}{\textbf{Model}} & \multicolumn{6}{c}{\textbf{File Level (\%)}} & \multicolumn{6}{c}{\textbf{Module Level (\%)}} & \multicolumn{6}{c}{\textbf{Function Level (\%)}} \\
        \cmidrule(lr){3-8} \cmidrule(lr){9-14} \cmidrule(lr){15-20}
        & & SR & REC & PRE & AJS & PLR & HR & SR & REC & PRE & AJS & PLR & HR & SR & REC & PRE & AJS & PLR & HR \\

        \midrule
        \multirow{2}{*}{SweRank}
        & Small & 8.00 & 14.84 & 11.13 & 6.64 & 0 & 31.11 & 4.17 & 9.44 & 5.24 & 3.45 & 0 & 24.07 & 1.02 & 2.85 & 1.28 & 0.74 & 0 & 11.73 \\
        & Large & 4.44 & 9.36 & 10.26 & 4.66 & 0 & 23.56 & 1.85 & 4.49 & 3.61 & 1.91 & 0 & 14.81 & 1.53 & 3.19 & 1.48 & 0.82 & 0 & 11.22 \\

        \midrule
        \multirow{2}{*}{Agentless}
        & GPT-4o & 13.33 & 18.83 & 15.66 & 10.07 & 2.22 & 32.44 & 8.33 & 12.86 & 6.93 & 4.23 & 0 & 24.54 & 5.10 & 8.02 & 2.04 & 1.49 & 0 & 14.29 \\
        & Claude-3.5 & 12.00 & 18.62 & 18.99 & 11.46 & 2.67 & 34.22 & 6.02 & 11.05 & 9.05 & 5.71 & 0.46 & 24.54 & 4.08 & 6.90 & 4.87 & 3.29 & 0 & 15.31 \\

        \midrule
        \multirow{2}{*}{LocAgent}
        & GPT-4o & 17.33 & 25.90 & 8.44 & 6.21 & 0.89 & 42.22 & 7.87 & 13.95 & 4.03 & 2.81 & 0 & 29.63 & 5.10 & 7.93 & 1.60 & 1.08 & 0 & 17.35 \\
        & Claude-3.5 & 22.22 & 36.78 & 11.02 & 8.73 & 0.44 & 54.67 & 15.74 & 26.85 & 6.63 & 4.97 & 0 & 47.69 & 11.73 & 20.79 & 4.32 & 3.22 & 0 & 40.82 \\

        \midrule
        \multirow{2}{*}{CoSIL}
        & GPT-4o & 10.22 & 17.12 & 15.62 & 9.39 & 2.22 & 32.89 & 6.94 & 11.69 & 9.47 & 5.06 & 0 & 24.54 & 3.06 & 4.98 & 3.42 & 1.87 & 0 & 10.71 \\
        & Claude-3.5 & 14.22 & 20.42 & 16.16 & 10.39 & 1.33 & 34.67 & 10.19 & 14.37 & 10.17 & 6.50 & 0.93 & 27.78 & 6.12 & 9.50 & 5.87 & 3.57 & 0 & 19.39 \\

        \midrule
        \multirow{2}{*}{\makecell{\tooldataloc}}
        & Claude-3.5 & 64.44 & 73.97 & 68.96 & 58.87 & 42.22 & \textbf{89.78} & 59.72 & 70.71 & 65.37 & 54.49 & 35.65 & \textbf{88.89} & 57.14 & 68.74 & 62.11 & 51.19 & 32.65 & \textbf{88.78} \\
        & Qwen3-Max & \textbf{68.44} & \textbf{77.74} & \textbf{73.35} & \textbf{65.28} & \textbf{48.44} & 87.11 & \textbf{62.96} & \textbf{75.19} & \textbf{70.40} & \textbf{61.44} & \textbf{42.59} & 86.11 & \textbf{58.67} & \textbf{72.55} & \textbf{66.98} & \textbf{58.03} & \textbf{38.27} & 84.69 \\

        \bottomrule
    \end{tabularx}
\end{table}
\end{landscape}



The granularity trend further exposes the limitations of existing localization methods under keyword-agnostic settings. Even under Hit Rate (HR), the most lenient metric that only requires at least one correct location, baselines degrade sharply from file-level to module-level and function-level localization. At the function level, even the best baseline results reach only 5.87\% precision and 3.57\% AJS. None of the baselines achieves a nonzero PLR. However, the strongest baseline recall is 20.79\%, indicating that partial retrieval remains possible even when exact localization is unsuccessful. In contrast, \tooldataloc achieves HR values ranging from 84.69\% to 89.78\% across both models and all three granularities. These results indicate that \tooldataloc maintains strong localization performance at finer granularities.

A key difference lies in how methods handle the size of the returned set. Baselines typically rely on top-$n$ recommendations to increase the probability of covering relevant locations. However, since the correct number of targets varies across queries and is not known in advance, this approach is inherently a compromise rather than an optimal solution. AJS and PLR assess the agreement between the returned sets and the ground truth. For instance, LocAgent with Claude-3.5 reaches a file-level HR of 54.67\%, but its AJS is only 8.73\%, suggesting that true positives are diluted within an inflated candidate set containing substantial noise. By contrast, \tooldataloc achieves substantially higher AJS across granularities, reflecting greater precision in returning constraint-satisfying locations. This precision is important for industrial deployment, as it reduces the validation burden for developers or downstream agents.

Abstention when no valid locations exist is further evaluated using \negset. The investigation of SWE-bench Lite shows that each issue instance involves only 1.15 ground-truth code changes on average. This sparsity raises a key question of whether existing tools truly pinpoint root causes or simply guess likely locations in a narrow search space. In \negset, query constraints are deliberately modified to create a mismatch between the query and the codebase, making the original ground-truth locations invalid. In this setting, the only correct behavior is to return \textit{no match found}. Unfortunately, all baselines suffer from a tendency to guess and persistently return top-$n$ recommendations even when constraints are not satisfied, yielding false positives. This may suggest that the strong performance of baselines is partly due to their recommendation-centric design rather than genuine localization reasoning. A lack of refusal ability can mislead downstream agents, waste computation, and introduce regression risk. In contrast, \tooldataloc returns a clear \textit{no match found} response for over 70\% of \negset queries, demonstrating the necessary discernment to abstain when ground-truth sets are empty.

\smallskip
\noindent\shadowbox{%
  \begin{minipage}{0.95\linewidth}
    \textbf{Answer to RQ1:} \tooldataloc outperforms the evaluated baselines on \dataset, attaining function-level precision of 66.98\% and a success rate of 58.67\%. It also returns \textit{no match found} for over 70\% of \negset queries.
  \end{minipage}}
